\chapter{Introduction}
\label{chap:intro}
%
\section{Background of the Study}

The global customer relationship management (CRM) software market was
valued at approximately USD~65.6~billion in 2023 and is projected to
reach USD~163.8~billion by 2030, representing a compound annual growth
rate (CAGR) of 13.9\% \citep{grandview2023}. This sustained growth
reflects the increasing digitisation of commercial operations, the
proliferation of customer interaction channels, and the rising
organisational demand for data-driven sales and retention strategies.

Traditional CRM implementations, which range from contact databases and
shared spreadsheets to legacy on-premise software suites, are not
equipped to handle the volume, velocity, or variety of modern customer
data. Nor can they surface the predictive intelligence that contemporary
sales organisations require to prioritise prospects and proactively retain
customers \citep{buttle2019}.

In Kazakhstan, the national digital economy agenda, formalised through the
Digital Kazakhstan state programme (2018--2022) and its successor Digital
Kazakhstan~2.0, has accelerated enterprise digitisation at a measurable
pace \citep{kazakhstan2024digital}. The number of registered
digital-commerce entities grew by approximately 47\% between 2020 and
2023, and more than 68\% of Kazakhstani SMEs reported using some form of
digital sales tool by 2024. Nevertheless, the adoption of sophisticated,
ML-augmented CRM platforms among domestic SMEs remains low. The majority
of organisations either rely on globally dominant incumbents such as
Salesforce or HubSpot, whose pricing structures are prohibitive in
tenge-denominated revenue environments, or maintain rudimentary spreadsheet
pipelines that preclude any form of predictive analysis.

Machine learning integration within CRM systems represents the principal
frontier of competitive differentiation in the sector. Predictive lead
scoring, defined as the assignment of a conversion probability to each
sales prospect, enables sales teams to concentrate effort on the
highest-value opportunities, materially improving funnel conversion
\citep{verbeke2012}. Customer churn analysis, that is, the identification
of at-risk customers before lapse occurs, supports proactive retention
interventions that the literature demonstrates can reduce churn rates by
20--30\% \citep{verbeke2012}. These capabilities, once accessible
exclusively to large enterprises with dedicated data science functions, are
increasingly realisable through mature open-source frameworks such as
scikit-learn \citep{chen2016xgboost} and interpretability toolkits such as
SHAP \citep{lundberg2017shap}.

\section{Problem Statement}

Despite the availability of globally recognised CRM platforms, Kazakhstani
SMEs face a compound set of structural barriers to adoption:

\begin{enumerate}

  \item \textbf{Cost barriers.} Salesforce CRM begins at USD~25 per user
  per month at the Essentials tier, while HubSpot's Sales Hub Professional
  is priced at USD~90 per user per month. These figures are prohibitive for
  most Kazakhstani SMEs operating within tenge-denominated revenue models
  and constrained foreign-currency budgets.

  \item \textbf{Data residency constraints.} Global CRM providers store
  customer data on servers located outside Kazakhstan, creating material
  compliance risk under the Law on Personal Data and its Protection
  (Law~No.~94-V, 2013, as amended 2022), which requires that personal data
  of Kazakhstani citizens be processed on servers physically situated
  within the Republic \citep{kazakhstanlaw94v}.

  \item \textbf{Absence of native ML intelligence.} Affordable CRM tools
  prevalent in the CIS market, such as AmoCRM and Bitrix24, provide basic
  pipeline management but offer no predictive lead scoring or churn
  prediction as native features. Obtaining these capabilities requires
  costly add-ons or bespoke integration work.

  \item \textbf{Localisation gaps.} The majority of global CRM platforms
  provide limited or no Kazakh-language support and are not optimised for
  locally dominant communication channels, such as WhatsApp, or for
  Kazakhstani payment gateway integration.

  \item \textbf{DevOps complexity.} Open-source self-hosted alternatives
  such as SuiteCRM and EspoCRM require significant infrastructure
  management expertise that exceeds the operational capacity of most SMEs
  in the market.

\end{enumerate}

The research problem addressed by this project is therefore:
\textit{How can a lightweight, ML-augmented CRM SaaS platform be designed
and prototyped that is technically feasible for a student development team,
cost-effective for Kazakhstani SMEs, compliant with domestic data
regulations, and deployable on Kazakhstani cloud infrastructure?}

\section{Purpose of the Project}

The purpose of this project is to design, develop, and evaluate a Minimum
Viable Product (MVP) of a SaaS CRM platform, designated \textit{TinyCRM},
that incorporates machine learning for predictive lead scoring and customer
churn analysis. The platform targets Kazakhstani SMEs and is deployed
exclusively on domestic cloud infrastructure. The project serves a dual
function: an academic purpose (demonstrating systems engineering and ML
integration competencies within a structured capstone framework) and a
practical purpose (producing a functional prototype suitable for future
commercial development).

\section{Objectives}

The project pursues the following specific objectives:

\begin{enumerate}[label=\textbf{O\arabic*.}]
  \item Conduct a systematic literature review of CRM architectures, SaaS
  design patterns, lead scoring methodologies, and churn prediction
  algorithms.

  \item Analyse five representative existing CRM systems to identify
  feature gaps relevant to the Kazakhstani SME context.

  \item Design a FastAPI/SQLAlchemy monolithic backend with a clearly
  specified RESTful API of 16 or more endpoints, supporting contact
  management, pipeline tracking, and ML inference.

  \item Develop and evaluate ML models for (a) lead scoring using Gradient
  Boosting, targeting F1 $\geq$ 0.78, and (b) churn prediction using
  Random Forest, targeting AUC-ROC $\geq$ 0.80.

  \item Design a Vue.js~3 single-page application providing a CRM
  dashboard, Kanban pipeline board, and ML insight panels.

  \item Deploy the MVP on QazCloud infrastructure and evaluate system
  performance under simulated concurrent load.

  \item Evaluate the system against functional and non-functional
  requirements and reflect critically on limitations and future work.
\end{enumerate}

\section{Scope and Limitations}

\subsection{Scope}

The project scope encompasses the following system boundaries:

\begin{itemize}
  \item \textbf{Backend:} A Python/FastAPI REST API with SQLAlchemy ORM,
  serving authentication, contact, company, deal, activity, and ML
  inference modules within a single deployable process.

  \item \textbf{Frontend:} A Vue.js~3 SPA utilising Pinia for state
  management and Vuetify~3 as the UI component library.

  \item \textbf{ML pipeline:} Data preprocessing, feature engineering,
  model training, serialisation via \texttt{joblib}, and REST-based
  real-time inference for lead scoring and churn prediction.

  \item \textbf{Database:} A PostgreSQL relational database with a
  defined schema supporting all CRM domain entities.

  \item \textbf{Deployment:} A Docker Compose single-server deployment on
  QazCloud virtual machines located in the Almaty region.

  \item \textbf{Testing:} Unit tests using \texttt{pytest}, API integration
  tests via Postman/Newman.
% , and manual user acceptance testing (UAT) sessions with three pilot SME participants.
\end{itemize}

